% Literature Review and Project Plan (MiSCADA summative submission).
% 8-page limit. Due 1 May 2026.
% Companion document: report/lit_review_standalone.pdf (full lit review, ~10 pages)
% which remains the complete reference treatment.

\documentclass[11pt,a4paper]{article}

\usepackage[utf8]{inputenc}
\usepackage[T1]{fontenc}
\usepackage{amsmath,amssymb,amsthm}
\usepackage[margin=2cm]{geometry}
\usepackage{hyperref}
\usepackage[longnamesfirst]{natbib}
\usepackage{booktabs}
\usepackage{array}
\usepackage{pgfgantt}
\usepackage{xcolor}
\bibliographystyle{plainnat}

\setlength{\parskip}{3pt}
\setlength{\parindent}{12pt}

\title{\vspace{-2em}Literature Review and Project Plan \\[2pt]
  \large Deep BSDE Methods for Multi-Agent Dealer Markets}
\author{Christian Garry \\ MSc MISCADA, Durham University}
\date{April 2026}

\begin{document}
\maketitle
\vspace{-1em}

\begin{abstract}
\noindent
The Cont--Xiong (2024) multi-agent dealer market defines a coupled
system of Hamilton--Jacobi--Bellman equations whose Nash equilibrium
is equivalent to the solution of a mean-field forward-backward
stochastic differential equation (FBSDE) with jumps.  This project
studies that FBSDE system computationally, using deep BSDE
methods as the numerical tool.  The aim is threefold: (i) to
approximate the $N$-player Nash equilibrium via a deep BSDE solver
and quantify its approximation error and convergence behaviour,
extending the curse-of-dimensionality-free theory of \citet{han2022}
to the jump-augmented setting; (ii) to characterise the conditions
under which naive deep-learning parameterisations fail to
represent the mean-field coupling of the HJB system, and to
design parameterisations that do; and (iii) to study the spread
gap $|S_{\mathrm{MADDPG}} - S_{\mathrm{Nash}}|$ as a quantitative
statistic for deviations of multi-agent reinforcement-learning
dynamics from the mean-field Nash equilibrium.  A sequential set of basic, intermediate, and
advanced goals, with an explicit test harness, structures the
13-week project timeline.
\end{abstract}

\vspace{-0.5em}

% =====================================================================
\section{Scientific Background}\label{sec:background}

Multi-agent dealer market-making sits at an uncomfortable intersection
in the mathematical literature.  The $N$-player stochastic-differential
game of \citet{cont2024} models competing market makers posting bid and
ask quotes in a continuous-time dealer market; executions arrive as
compensated point processes, and Nash equilibrium is characterised by
a coupled system of Hamilton--Jacobi--Bellman (HJB) equations.  The
game admits \emph{in principle} a probabilistic treatment via backward
stochastic differential equations (BSDEs) and mean-field games (MFGs),
and a numerical treatment via deep-learning methods.  \emph{In
practice}, no existing framework simultaneously handles all three
ingredients that the model requires:
\citet{han2022} give a curse-of-dimensionality-free deep-BSDE method
for McKean--Vlasov forward-backward SDEs but without jumps;
\citet{wang2023} extend deep BSDE to jump-diffusion processes but
without mean-field coupling; and \citet{cont2024} solve their own
game empirically via multi-agent deep deterministic policy gradient
(MADDPG) but without the equilibrium grounding or interpretable
gradients that a principled method would supply.  Combining all three
---~jump-augmented BSDEs, mean-field coupling, and deep-learning
numerics~--- into a single computational method for a realistic market
microstructure model is the gap this project addresses.

Concretely, each dealer controls the bid/ask spread $\delta_t$ over
their quote, and the execution probability in response to a random
order is governed by the softmax-type kernel of \citet{cont2024}
(their eq.~58),
\begin{equation}\label{eq:cx-exec}
  f_a^i(\delta, s) \;=\;
  \tfrac{1}{N}\,\sigma(-\delta)\,
  \frac{e^{s}}{1 + e^{\,\delta + s}},
\end{equation}
where $\sigma(\cdot)$ is the logistic sigmoid and $s = S/\kappa$ is
the market-price scale.  Nash equilibrium is the fixed point of the
coupled Bellman system, and existence follows from a Brouwer
fixed-point argument; what is missing is a scalable numerical
method that respects both the mean-field limit and the
jump-diffusion structure.


% =====================================================================
\section{Related Work}\label{sec:related}

The literature divides naturally into three pillars.

\paragraph{BSDEs and mean-field games.}
\citet{pardoux1990} proved existence and uniqueness of adapted BSDE
solutions under Lipschitz hypotheses, providing the bedrock on which
\citet{elkaroui1997} built the BSDE formulation of derivative pricing
and risk management.  \citet{feng2018} extended the PDE/FBSDE
correspondence to the \emph{weak}-solution setting for quasi-linear
PDEs --- directly relevant to the Cont--Xiong HJB system, whose
coefficients are quasi-linear rather than classical.  A BSDE on $[0,T]$ is a pair of adapted
processes $(Y_t, Z_t)$ satisfying
\begin{equation}\label{eq:bsde}
  Y_t \;=\; \xi + \int_t^T f(s, Y_s, Z_s)\,ds - \int_t^T Z_s \,dW_s,
\end{equation}
with $\xi$ a terminal condition; the $Z$-process is the object that
deep BSDE methods learn.  \citet{peng2004} established the link
between BSDEs and nonlinear expectations, placing $Z$ in the role
of a risk-sensitivity functional.  Mean-field games were introduced
by Lasry--Lions and Huang--Caines--Malham\'e; \citet{carmona2013}
and the two-volume monograph of \citet{carmona2018} supply the
probabilistic treatment via McKean--Vlasov FBSDEs, and
\citet{buckdahn2009viscosity, buckdahn2009limit} establish the
MF-BSDE / non-local PDE equivalence and the $1/\sqrt{N}$ propagation
of chaos rate that justifies mean-field approximation at finite
$N$.  The jump-augmented case has a parallel classical theory
\citep{tang1994, barles1997} on which any deep BSDEJ method must
rest.

\paragraph{Deep BSDE numerics.}
The curse of dimensionality in classical BSDE schemes
\citep{kloeden1992, pengxu2011} --- which require
conditional-expectation estimation scaling as $O(M^d)$ in state
dimension --- was broken by the deep BSDE method of \citet{han2018},
which parameterises the gradient $Z_t$ by a neural network and trains
on the terminal-loss $|Y_T - \xi|^2$.  \citet{beck2019} extended
this to fully nonlinear second-order BSDEs.  For mean-field problems,
\citet{han2022} proved curse-of-dimensionality-free convergence for
McKean--Vlasov FBSDEs via a fictitious-play decomposition
\citep{guo2023}, with the caveat that the proof requires Lasry--Lions
monotonicity on the coupling term --- a condition that for the
Cont--Xiong kernel~\eqref{eq:cx-exec} can only be verified empirically
along a one-dimensional slice of the measure space.  The extension
to jump-diffusion processes and forward BSDEs with jumps (FBSDEJs) is
due to \citet{wang2023}; the BSDEJ numerical literature is as yet
thinly cited, and any downstream application must do its own due
diligence on convergence and on the compensated-martingale correction.

\paragraph{Structural pathologies of deep approximations to the
MFG HJB system.}
\citet{germain2022} catalogue vanishing gradients and terminal-scale
sensitivity for the canonical deep BSDE.  Three further failure
modes are specific to the mean-field and jump-augmented setting
and are structural --- each corresponds to a specific way in
which the neural-network approximation fails to represent the
mean-field HJB system it is supposed to solve, rather than a
generic optimisation difficulty.

\emph{(i) Loss of measurable $\mu$-dependence via BatchNorm.}
For a mean-field game the value function is $V(t, x, \mu_t)$ and
the HJB equation it satisfies is explicitly $\mu$-dependent.  When a
broadcast population statistic $m_t$ (a low-dimensional surrogate
for $\mu_t$) is included as a network input and passed through a
BatchNorm layer, normalisation across the batch dimension maps
the constant feature to zero, so the parameterised value function
becomes measurable with respect to own-state only.  The
discretised system effectively solves a single-agent HJB, not
the MFG.

\emph{(ii) Degenerate fixed points of symmetric-input DeepSets.}
Permutation-invariant pooling of opponent states is the
architecturally correct way to encode the empirical distribution.
Under symmetric initialisation and symmetric inputs, however,
DeepSets admits degenerate minima in which the map from
population state to policy collapses to a constant.  The effect
is operationally indistinguishable from (i): the approximation
again loses measurable dependence on $\mu$.

\emph{(iii) Underdetermined terminal loss (``generator bypass'').}
The deep BSDE terminal loss $|Y_T - g(q_T)|^2$ is a single scalar
constraint on the network-parameterised $U$-process.  Many
$U$-profiles satisfy it in expectation, but only one corresponds
to the Bellman-optimal $U^\star(q) = V(q-1) - V(q)$ of the
underlying HJB system.  Without a prior, Adam selects whichever
basin is numerically easiest, and the method can converge to a
self-consistent fixed point of the BSDE discretisation that does
not solve the HJB --- the spread profile extracted from such a
$U$ is not a best response to the population quote.

Diagnosing and repairing these three pathologies is a core
methodological contribution of the project.

\paragraph{Multi-agent RL and tacit collusion.}
\citet{cont2024} show via decentralised MADDPG that independently
learning market-making algorithms converge to spreads strictly above
the Nash equilibrium, a phenomenon termed \emph{tacit collusion}.
Related MARL work on realistic limit-order-book simulators
\citep{karpe2020, shi2024, cheridito2025} highlights fragility across
market regimes.  The gap between MADDPG spreads and Nash,
$|S_{\mathrm{MADDPG}} - S_{\mathrm{Nash}}|$, is a measurable
signature of algorithmic collusion, but it has not been operationalised
as a detector statistic --- which requires a BSDE-based Nash benchmark
that is itself trusted.  Recent frontier work on common-noise
MV-FBSDEs via path signatures \citep{hu2025} and on mean-field LOB
embeddings \citep{navarro2025} defines the open problems beyond this
project's scope.


% =====================================================================
\section{Research Question}\label{sec:question}

\emph{What are the approximation error and convergence
behaviour of a deep BSDE-with-jumps approximation to the
$N$-player Nash equilibrium of the Cont--Xiong mean-field FBSDE,
and what structural conditions on the neural parameterisation
preserve the mean-field HJB coupling under that approximation?}
Three sub-questions structure the investigation.  (Q1)~What is
the approximation error $|S^{\text{neural}}_0 - S^{\text{Nash}}_0|$,
and its scaling in $N$, for a jump-augmented deep BSDE solver
applied to this model; in particular, is the error consistent
with the $O(1/\sqrt{N})$ propagation-of-chaos rate of
\citet{buckdahn2009limit}?  (Q2)~Which structural pathologies of
the neural parameterisation (\S\ref{sec:related}) admit numerical
fixed points that do not solve the underlying HJB system, and
which architectural or regularisation choices eliminate them?
(Q3)~How does the spread gap
$|S_{\mathrm{MADDPG}} - S_{\mathrm{Nash}}|$ --- a quantitative
statistic for deviation of multi-agent RL dynamics from the
mean-field Nash equilibrium --- behave across model parameters
$(N, \lambda, \phi)$?


% =====================================================================
\section{Anticipated Outcomes}\label{sec:outcomes}

Outcomes are stratified into three levels.  Tier~1 (basic) is
essentially complete and de-risks the submission.  Tier~2
(intermediate) is in progress with known technical challenges.
Tier~3 (advanced) is genuinely speculative and is the ceiling,
not the floor.

\begin{table}[h!]
\centering
\small
\begin{tabular}{>{\raggedright\arraybackslash}p{0.22\textwidth}
                >{\raggedright\arraybackslash}p{0.62\textwidth}
                >{\centering\arraybackslash}p{0.1\textwidth}}
\toprule
\textbf{Tier} & \textbf{Outcome} & \textbf{Risk} \\
\midrule
T1: Basic
  & Exact $N$-player Nash equilibrium of Cont--Xiong via policy-evaluation
    fictitious play on a discrete inventory grid, validated at
    $N = 2, 5, 10, 25, 100$ against the Bellman fixed-point tolerance.
    Shared-weight finite-horizon BSDEJ solver matching Nash spread at $q=0$
    to within 1\% with fixed $\mathrm{avg\_comp}$.
  & \textcolor{green!50!black}{Low} \\[2pt]
T1: Basic
  & Pytest-level unit tests and regression tests for all solvers
    (boundary-quote convention, compensated-martingale correction,
    FOC optimality, propagation-of-chaos sanity).
  & \textcolor{green!50!black}{Low} \\[2pt]
\midrule
T2: Intermediate
  & Mean-field fictitious-play BSDEJ solver converging to Nash at
    $N = 2, 5, 10$ with spread error $< 3\%$.  Requires resolving the
    generator-bypass failure mode identified in \S\ref{sec:related}.
  & \textcolor{orange}{Medium} \\[2pt]
T2: Intermediate
  & Empirical verification of $O(1/\sqrt{N})$ propagation of chaos
    \citep{buckdahn2009limit} at $N \in \{2, 3, 5, 10, 25, 100, 5000\}$
    via exact solver, and tightness of the bound for the neural solver.
  & \textcolor{orange}{Medium} \\[2pt]
T2: Intermediate
  & Architecture ablation study (ReLU vs.\ Tanh vs.\ GELU; shared vs.\
    per-timestep networks; DeepSets vs.\ mean-pooling) quantifying each
    failure mode in \S\ref{sec:related}.
  & \textcolor{green!50!black}{Low} \\[2pt]
\midrule
T3: Advanced
  & Continuous-inventory diffusion BSDE solver with $Z \neq 0$,
    learning $Z_t = \sigma_q\,\partial_q V$ genuinely rather than the
    pure-jump degenerate case.
  & \textcolor{red!70!black}{High} \\[2pt]
T3: Advanced
  & MADDPG reproduction of \citet{cont2024} and quantitative
    operationalisation of the spread gap
    $|S_{\mathrm{MADDPG}} - S_{\mathrm{Nash}}|$ as a collusion
    detector across a parameter sweep in $(N, \lambda, \phi)$.
  & \textcolor{red!70!black}{High} \\[2pt]
T3: Advanced
  & Common-noise extension via path-signature conditioning
    \citep{hu2025}; multi-asset extension ($K = 2, 3, 5$) demonstrating
    genuine curse-of-dimensionality benefit over the exact solver.
  & \textcolor{red!70!black}{High} \\
\bottomrule
\end{tabular}
\end{table}

Tier~1 is essentially complete and provides a fall-back position
should the more ambitious components prove intractable.  Tier~2
contains the principal research content of the project.  Tier~3
represents original extensions that may or may not be achievable
within the available timeframe; they are reported in the
dissertation if completed, and identified as clearly defined open
problems if not.


% =====================================================================
\section{Project Time Plan}\label{sec:plan}

The project runs from project allocation (January 2026, completed)
through submission on 4 September 2026.  The core dissertation phase
runs 27 April -- 26 June 2026, with July--August reserved for
write-up.  A Gantt summary of the remaining timeline is below.

\begin{center}
\begin{ganttchart}[
    x unit=0.55cm,
    y unit title=0.65cm,
    y unit chart=0.55cm,
    hgrid, vgrid,
    title label font=\footnotesize,
    bar label font=\scriptsize,
    milestone label font=\scriptsize,
    group label font=\scriptsize\bfseries,
    title height=1,
    bar height=.6,
    group height=.6
]{1}{19}
\gantttitle{April}{2}
\gantttitle{May}{4}
\gantttitle{June}{5}
\gantttitle{July}{4}
\gantttitle{Aug}{4}\\
\gantttitlelist{1,...,19}{1}\\
\ganttgroup{T1 (done)}{1}{1} \\
\ganttbar{Exact Nash + BSDEJ baseline}{1}{1} \\
\ganttgroup{Submission \& feedback}{1}{3} \\
\ganttmilestone{Lit review \& plan due}{2} \\
\ganttbar{Supervisor feedback}{2}{3} \\
\ganttgroup{T2: MV-FP BSDEJ}{3}{8} \\
\ganttbar{Fix generator bypass}{3}{5} \\
\ganttbar{$N=2,5,10$ validation}{5}{7} \\
\ganttbar{$1/\sqrt{N}$ rate verification}{6}{8} \\
\ganttbar{Architecture ablations}{7}{9} \\
\ganttgroup{T3: Advanced}{9}{13} \\
\ganttbar{Continuous-$q$ diffusion BSDE}{9}{11} \\
\ganttbar{MADDPG + collusion detector}{10}{13} \\
\ganttbar{Multi-asset $K$-scaling}{11}{13} \\
\ganttmilestone{Mid-May demo}{4} \\
\ganttgroup{Write-up}{13}{18} \\
\ganttbar{Method + results chapters}{13}{16} \\
\ganttbar{Intro + conclusion + revisions}{15}{18} \\
\ganttmilestone{Draft to supervisor}{17} \\
\ganttbar{Final polish}{18}{19} \\
\ganttmilestone{Submission}{19}
\end{ganttchart}
\end{center}

Weekly numerals on the chart correspond to one-week periods with
week~1 being 27~April 2026 (start of the core dissertation phase).
Weekly supervisor meetings are scheduled throughout Term~3; the
preliminary demonstration (MiSCADA milestone) is pencilled in for
mid-May.  The draft handover to the supervisor is scheduled for
the end of the third week of July to allow a round of feedback
before the September deadline.


% =====================================================================
\section{Demonstrating Success}\label{sec:success}

The project is a numerical one and admits quantitative validation at
every tier.

\paragraph{Ground-truth validation.}
The classical fictitious-play Bellman solver provides an exact
reference Nash equilibrium for discrete inventory and small $N$.
All neural solvers are validated against this reference.  The
principal success metric is the \emph{relative spread error at
$q=0$},
\begin{equation*}
  \varepsilon_{\text{spread}} \;=\; \frac{|\,S^{\text{neural}}_0 - S^{\text{Nash}}_0\,|}{S^{\text{Nash}}_0},
\end{equation*}
with targets of $\varepsilon_{\text{spread}} < 1\%$ at Tier~1 (fixed
population quote) and $\varepsilon_{\text{spread}} < 3\%$ at Tier~2
(MV-fictitious-play closure).  The current best Tier~1 result is
0.59\% error at $N = 2$.

\paragraph{Empirical convergence-rate tests.}
The $O(1/\sqrt{N})$ propagation-of-chaos bound of
\citet{buckdahn2009limit} gives a principled, non-asymptotic check:
computing the exact Nash at $N \in \{2, 3, 5, 10, 25, 100, 5000\}$ and
fitting $S^{\text{Nash}}_0(N) = a + b/\sqrt{N}$ should yield a
residual consistent with the theoretical rate.  Deviations flag
either a solver bug or a breakdown of the mean-field regime.

\paragraph{Failure-mode diagnostics.}
Each of the three failure modes identified in \S\ref{sec:related}
admits a \emph{targeted} test.  BatchNorm erasure is detected by
injecting a known $m_t$-dependent signal and verifying that the
network's output responds to it (controlled-input test).
Representation collapse under DeepSets is detected by permuting
opponent states and measuring the deviation of the network output
(equivariance test).  Generator bypass is detected by comparing the
network's extracted optimal quotes against the FOC solution to the
analytic optimisation problem at each inventory state
(FOC-consistency test).

\paragraph{Reproducibility and regression.}
All solvers are covered by a pytest suite
(\texttt{tests/test\_*.py}) pinned to small problem sizes
($N = 2$, $Q = 5$, $M = 50$ time steps) and numerical tolerances
that detect any regression in the compensated-martingale correction
or boundary-quote convention.  Intermediate results are saved as
JSON under \texttt{results\_final/} for incremental recovery on
failure, and all experiments are seeded for deterministic replay.

\paragraph{Demonstrating negative results.}
The regulations explicitly note that negative results are still
results.  The project anticipates at least one such outcome:
characterising precisely \emph{why} the naive MV-fictitious-play
BSDEJ converges to spurious self-consistent fixed points (as
observed in early prototyping) is itself a contribution, whether
or not the repair succeeds.  Partial success is reported
transparently, with failure modes documented rather than papered
over.


% =====================================================================
\small
\bibliography{dissertation_study/references}

\end{document}
